\documentclass[10pt,a4paper]{amsart}
\usepackage[margin=19mm]{geometry}
\usepackage{amsmath,amssymb,mathtools}
\usepackage[scr=rsfs,cal=euler]{mathalfa}
\usepackage{xfrac}
\usepackage[unicode,hidelinks,bookmarks=false]{hyperref}
\newcommand{\ie}{i.e.\ }
\newcommand{\cf}{cf.\ }
\newcommand{\resp}{respectively\ }
\newcommand{\Subsection}{\subsection}
\providecommand{\nobreakdash}{\nobreak\hbox{-}}
\setlength{\emergencystretch}{2em}
\multlinegap0pt
\usepackage{siunitx}
\usepackage{siunitx}
\usepackage{siunitx}
\usepackage{siunitx}
\usepackage{amssymb}
\usepackage{float}
\usepackage{xcolor}
\usepackage[shortlabels]{enumitem}
\usepackage{mathtools}
\usepackage{siunitx}
\usepackage{algorithm}\usepackage{algpseudocode}
\usepackage[normalem]{ulem}
\newtheorem{theorem}{Theorem}
\title{Tensor-based empirical interpolation method,\newline and its application in model reduction}
\author{Brij Nandan Tripathi, Hanumant Singh Shekhawat, Seip Weiland}
\date{Source: arXiv:2410.21770v2 (CC BY 4.0)}
\begin{document}
\maketitle

\noindent\emph{Source and licence.} Tensor-based empirical interpolation method,\newline and its application in model reduction. Version arXiv:2410.21770v2 (CC BY 4.0), distributed by arXiv under the Creative Commons Attribution 4.0 International licence, http://creativecommons.org/licenses/by/4.0/, as recorded in the arXiv licence metadata for this exact version and shown on the abstract page https://arxiv.org/abs/2410.21770v2. Full licence notice: \texttt{licenses/arxiv-2410.21770-CC-BY-4.0.txt}.\par\smallskip

\noindent\textit{Editorial scope.} Counted unit: the article's theorem lem63 (source0.tex lines 651--653). The article's complete original proof of it is delivered with it (source0.tex lines 654--681, one \texttt{proof} environment of 28 lines). No mathematics is written, repaired or re-derived: the statement and the proof are the article's own lines, in the article's own notation, and no retained line is substituted or re-flowed.  Retained uncounted as the source-local context the proof needs: lines 581--586.  Nothing was omitted from any retained span, so the package carries no omission record.  No retained line contains a citation, so the wrapper carries no bibliography and no bibliography entry.  No retained line contains a figure environment, an image-inclusion command or drawing code, so no image is delivered and the package declares no asset dependency.  Editorial formatting only, in addition: an AMS \texttt{amsart} preamble that restates the article's own declarations and macros used by the retained lines, namely the environments \texttt{theorem} on separate counters, together with the standard packages those lines need.  The retained environments print in this document as Theorem 1 in order of appearance, whereas the article prints the same items with the numbering of its own full document.  The complete native source of the article as distributed in the arXiv e-print for this exact version is bundled unchanged as \texttt{sources/167569/source0.tex}.
\begin{equation}\label{eq511}
     \begin{split}
         R_{(k-1)m_2+l}&=(\bold{u_k}-\sum_{j=1}^{k-1}\bold{c}_{\bold{((j-1)m_2+\ell},(k-1)m_2+1)}\bold{u_j})\\
                        &\quad\quad\quad*(\bold{v_\ell}^T-\sum_{i=1}^{\ell-1}\bold{c}_{(\bold{\ell},i)}\bold{v_i}^T)
     \end{split}
 \end{equation}

\begin{theorem}\label{lem63}
Suppose $P_5$ and $P_6$ are the masking operators obtained by applying the Discrete Empirical Interpolation Algorithm on columns of $U_1$ and $U_2$, respectively. Then $P_3=P_5$ and $P_4=P_6$.
\end{theorem}

\begin{proof}
    As per equation \eqref{eq511},  the residual conforms to the following structure for any given iteration set $s_k$. Within this iteration set $s_k$, the same structure holds true for any arbitrary element $l$.
\begin{equation*}
     \begin{split}
        max(\lvert R_{(k-1)m_2+l} \rvert)&=max(\lvert \bold{u_k}-\sum_{j=1}^{k-1}\bold{c}_{\bold{((j-1)m_2+\ell},(k-1)m_2+1)}\bold{u_j} \rvert)\\
        & \quad \quad *max(\lvert\bold{v_\ell}^T-\sum_{i=1}^{\ell-1}\bold{c}_{(\bold{\ell},i)}\bold{v_i}^T \rvert)\\
                            &= max(r_1)max(r_2)
     \end{split}
 \end{equation*}
% Where $r_1= \bold{u_k}-&\sum_{j=1}^{k-1}\bold{c}_{(\bold{(j-1)m_2+\ell},(k-1)m_2+1)}\bold{u_j}$ and $r_2=\bold{v_\ell}^T-\sum_{i=1}^{\ell-1}\bold{c}_{(\bold{\ell},i)}\bold{v_i}^T$. 
\begin{equation*}
\text{Where } r_1 = \mathbf{u_k} - \sum_{j=1}^{k-1} \mathbf{c}_{\left( (j-1)m_2 + \ell , (k-1)m_2 + 1 \right)} \mathbf{u_j} 
\text{ and } r_2 = \mathbf{v_\ell}^T - \sum_{i=1}^{\ell-1} \mathbf{c}_{(\ell, i)} \mathbf{v_i}^T
\end{equation*}


For any arbitrary iteration k and $\ell$, residual is of the below form in the DEIM algorithm for the columns of $U_1$ and $U_2$, respectively 
  \begin{equation*}
     \begin{split}
         R_{k}&=(\bold{u_k}-\sum_{j=1}^{k-1}\bold{c}_{(\bold{k},j)}\bold{u_j})
     \end{split}
 \end{equation*}
 \begin{equation*}
     \begin{split}
         R_{l}&=(\bold{v_\ell}-\sum_{j=1}^{\ell-1}\bold{c}_{(\bold{\ell},j)}\bold{v_j})
     \end{split}
 \end{equation*}
\end{proof}

\end{document}
